\section{Алгоритмы информационно-поисковой системы}
\label{sec:algorithms}

Разработанная система реализует вероятностную стратегию поиска по варианту~5 задания: документы ранжируются по убыванию оценки вероятности релевантности запросу. Пользователь формулирует запрос на естественном русском языке; система автоматически строит поисковый образ запроса, сопоставляет его с инвертированным индексом коллекции и возвращает упорядоченный список документов.

\subsection{Предобработка текста}

Документы и запросы проходят одинаковый конвейер преобразований, чтобы поисковые образы имели одинаковую структуру:

\begin{enumerate}
	\item приведение к нижнему регистру и замена буквы <<ё>> на <<е>>;
	\item токенизация регулярным выражением, допускающим внутренние дефисы;
	\item лемматизация русскоязычных словоформ анализатором \texttt{pymorphy3};
	\item фильтрация стоп-слов и токенов короче двух символов.
\end{enumerate}

Латинские термины и аббревиатуры (SQL, IP, Ethernet) сохраняются в исходном написании: морфологический словарь русского языка к ним неприменим, а сами строки уже являются нормальной формой. Результат предобработки "--- список лемм, образующих поисковый образ документа или запроса.

\subsection{Инвертированный индекс}

Для каждого документа строится инвертированный индекс: термин указывает на список документов, в которых он встречается, с указанием частоты. Дополнительно хранятся документная частота термина $n_i$ и длина документа $|d|$ в леммах. Эти статистики достаточны для вычисления весов Робертсона -- Спарк Джонс и оценки BM25 без повторного чтения текстов коллекции.

\subsection{Вероятностная модель ранжирования}

В основе лежит принцип вероятностного ранжирования: оптимальная выдача упорядочивает документы по убыванию вероятности их релевантности запросу. На практике вычисляется монотонное преобразование этой вероятности "--- статус релевантности RSV.

Вес термина без обратной связи задаётся сглаженной формулой Робертсона -- Спарк Джонс:

\begin{equation}
	w_i = \log\left(1 + \frac{N - n_i + 0{,}5}{n_i + 0{,}5}\right),
	\label{eq:rsj}
\end{equation}

где $N$ "--- число документов в коллекции; $n_i$ "--- число документов, содержащих термин $i$. Слагаемое $1$ внутри логарифма устраняет отрицательные веса для терминов, встречающихся более чем в половине документов, сохраняя монотонность: редкий термин получает больший вес, чем частый.

Ранжирование выполняется по формуле Okapi BM25:

\begin{equation}
	\operatorname{RSV}(d,q) = \sum_{i \in q} w_i \cdot
	\frac{f_{i,d}\,(k_1 + 1)}{f_{i,d} + k_1\left(1 - b + b\,\dfrac{|d|}{\operatorname{avgdl}}\right)},
	\label{eq:bm25}
\end{equation}

где $f_{i,d}$ "--- частота термина $i$ в документе $d$; $|d|$ "--- длина документа; $\operatorname{avgdl}$ "--- средняя длина документа коллекции. Параметр $k_1 = 1{,}2$ задаёт насыщение частоты: вклад термина растёт быстро при первых вхождениях и выходит на плато. Параметр $b = 0{,}75$ нормирует оценку по длине документа, не давая преимущество длинным текстам только за счёт объёма.

\subsection{Псевдорелевантная обратная связь}

Полная вероятностная модель предполагает известное множество релевантных документов. Пользователь его не задаёт, поэтому применяется псевдорелевантная обратная связь. После первого прохода BM25 верхние $R = 5$ документов объявляются условно релевантными, и вес термина пересчитывается:

\begin{equation}
	w_i = \log\frac{(r_i + 0{,}5)\,/\,(R - r_i + 0{,}5)}
	{(n_i - r_i + 0{,}5)\,/\,(N - n_i - R + r_i + 0{,}5)},
	\label{eq:rsj-prf}
\end{equation}

где $R$ "--- размер псевдорелевантного множества; $r_i$ "--- число документов этого множества, содержащих термин $i$. Второй проход BM25 выполняется уже с уточнёнными весами. В интерфейсе обратная связь включается отдельным флажком: на малой разнородной коллекции она может понижать точность, подмешивая общие термины из соседних документов.

\subsection{Отсечение нерелевантного хвоста}

Оценка BM25 положительна для любого документа, содержащего хотя бы один термин запроса. Чтобы выдача не засорялась слабо связанными текстами, документы, чья оценка ниже $0{,}25$ от оценки лидера, отбрасываются. Порог подобран экспериментально: он сохраняет тематически близкие результаты и отсекает случайные совпадения.

\subsection{Базовая векторная модель}

Для сравнения реализована векторная модель с логарифмическим TF-IDF и косинусной близостью:

\begin{equation}
	\operatorname{sim}(d,q) = \frac{(\vec{D},\vec{Q})}{\lVert\vec{D}\rVert\cdot\lVert\vec{Q}\rVert},
	\qquad
	w^{d}_{i} = \bigl(1 + \log f_{i,d}\bigr)\cdot\bigl(\log(N/n_i) + 1\bigr).
	\label{eq:tfidf}
\end{equation}

Эта модель не используется как основная: она соответствует другому варианту задания и служит только контрольной точкой при оценке качества.

\subsection{Метрики качества поиска}

Пусть $Rel$ "--- множество эталонных релевантных документов по запросу, $R$ "--- множество документов, выданных системой. Тогда

\begin{equation}
	P = \frac{|R \cap Rel|}{|R|},
	\qquad
	\operatorname{Recall} = \frac{|R \cap Rel|}{|Rel|},
	\qquad
	F_{\beta} = \frac{(1 + \beta^{2})\,P\cdot\operatorname{Recall}}{\beta^{2}\,P + \operatorname{Recall}}.
	\label{eq:prf}
\end{equation}

Дополнительно вычисляются точность на глубине $k$ ($P@5$, $P@10$), R-Precision, средняя точность AP, усреднённая по запросам MAP, нормализованный дисконтированный выигрыш nDCG@10 и 11-точечная интерполированная кривая полноты-точности. Параметр $\beta$ задаёт акцент: $F_{0{,}5}$ сильнее штрафует падение точности, $F_2$ "--- падение полноты, $F_1$ уравновешивает оба свойства.
